\section{Robustness}
\label{sec:robustness}

We address the two most demanding threats---the 2023 banking stress and
inference under a common adoption quarter---then pre-trends, the scope
of the response, and a set of compressed additional checks.

\subsection{The 2023 Stress: Characteristics-by-Quarter Controls}
\label{sec:robustness:stress}

If high-CECL banks were differentially exposed to the March 2023 stress
through observable balance-sheet characteristics, quarter fixed effects
alone would not remove the confound.  We therefore allow the eight
2022Q4 characteristics the literature identifies as run-risk
drivers---log assets, equity-to-assets, uninsured deposit share,
brokered share, NPL/loans, ROA, CRE share, and unrealized securities
losses \citep{caglio2024depositors}---to enter interacted with
calendar-quarter fixed effects.  In the level DiD, size and capital
interactions leave the estimate essentially unchanged ($-0.069$,
$p<0.10$), while adding quarter-specific slopes in the 2022Q4 deposit-mix
variables attenuates it ($-0.037$ continuous, $-0.238$ binary,
unreported table).  Those deposit-mix controls are mechanically
entangled with the outcome---the bad-control problem that motivates the
composition design.  Table~\ref{tab:flexcontrols_td} applies the same
saturation \emph{within} the triple difference, where the entanglement
does not arise; the triple interaction survives essentially intact
($-0.136$ per percentage point of exposure, $p<0.05$; $-1.05$ percentage
points for high-CECL banks, $p<0.01$; baseline: $-0.177$ and $-1.25$).
Whatever differential stress exposure the 2022Q4 characteristics carry,
it does not explain why uninsured deposits fell \emph{relative to
insured deposits within the same bank} in proportion to the disclosed
adjustment.  The weekly posted-rate evidence of
Section~\ref{sec:results:anticipation} reinforces the point from
outside the Call Report: high-CECL banks showed no differential retail
funding-market response during the stress window itself.

\subsection{Inference under a Common Adoption Quarter}
\label{sec:robustness:inference}

Because 92 percent of the sample adopts CECL in 2023Q1, the design has a
strong common-shock component, and standard errors clustered only at the
bank level could understate uncertainty if cross-sectional residuals
share the post-shock disturbance.  Table~\ref{tab:inference_summary}
addresses this two ways.  First, two-way clustering by bank and calendar
quarter leaves the standard errors of the headline estimates essentially
unchanged (the triple-difference standard error in fact falls slightly).
Second, we conduct randomization inference in the spirit of a
design-based Fisher test: we reassign the bank-level treatment pair
(continuous adjustment and the derived High-CECL indicator) at random
across banks, holding each bank's adoption timing fixed, re-estimate the
headline specifications on 1{,}000 permuted assignments, and compute
two-sided $p$-values as the share of placebo coefficients at least as
large in magnitude as the actual estimate.  The composition
triple-difference has a permutation $p$-value of 0.001 in both
continuous and binary form; interest expense (0.006) and ROA (0.001)
are similarly sharp.  The level uninsured DiD is marginal under
permutation ($p \approx 0.05$), consistent with our anchoring the paper
on the composition contrast rather than the level result.

\subsection{Cohort Timing and the Not-Yet-Adopted Placebo}
\label{sec:robustness:timing}

The staggered rollout offers a lever in principle: 323 banks adopted in
quarters other than 2023Q1.  In practice this margin has limited power,
and we do not lean on it.  Dropping the entire 2023Q1 cohort leaves only
13 high-CECL banks; the resulting DiD carries the right sign (continuous
$-0.081$; binary $-0.133$) but is imprecise.

A placebo on the 231 late-cohort banks that had not yet adopted as of
2023Q1---regressing the 2023Q1 change in uninsured time deposits on the
size of their still-undisclosed \emph{future} Day-One adjustment---does
not come out clean: the coefficient is negative and significant
($-0.101$ continuous, $p<0.05$).  Two readings are consistent with this
pattern.  Under the first, the March 2023 stress loaded on the same
credit-risk fundamentals that CECL later capitalized into a single
number, and the placebo captures that common fundamental.  Under the
second, forward-looking depositors anticipated the disclosure itself in
the weeks between the March failures and the 2023Q1 filings.  We do not
view the placebo as decisive between these readings.  What it
establishes is that the level DiD around 2023Q1 cannot by itself
separate disclosure from stress at the primary cohort---which is
precisely why identification rests on the characteristics-by-quarter
saturation of the triple difference and on the private-versus-public
timing evidence of Section~\ref{sec:results:anticipation}, neither of
which this cross-sectional 2023Q1 correlation can contaminate.

\subsection{Extended Pre-Trends, Including the Pandemic Quarters}
\label{sec:robustness:pretrends}

The baseline event window begins in 2020Q3.
Figure~\ref{fig:es_long_balanced} extends the event study back to 2018Q1
with endpoint bins at $k=-20$ and $k=+11$, so that the thinly populated
edge quarters contributed only by off-cycle cohorts are pooled rather
than estimated from a handful of banks.  Point estimates for uninsured
time deposits over the five-year pre-period sit near zero with no slope
toward the adoption-date divergence, and the post-adoption break remains
the dominant feature.  Joint tests of the pre-adoption coefficients do
not reject parallel trends (uninsured $p = 0.18$; insured $p = 0.37$).
In earlier unbinned versions the same tests rejected marginally
($p \approx 0.04$--$0.06$); the rejections were driven entirely by the
unbalanced edge quarters, not by any trend near adoption.  The COVID
quarters contribute noise but no differential trend.  The one visible
pre-adoption movement---the late upward drift in the insured
series---is the balance-sheet footprint of the bank-side pre-funding
documented in Section~\ref{sec:results:anticipation}.

\subsection{Scope of the Response: Total Uninsured Deposits}
\label{sec:robustness:scope}

Table~\ref{tab:total_uninsured} estimates the baseline DiD on total
uninsured and total insured deposits scaled by assets.  Neither moves
detectably (total uninsured: continuous $+0.073$, SE $0.085$; binary
$-0.177$, SE $0.458$).  The disciplinary response we document is
confined to the time-deposit rollover margin.  We regard this as the
mechanism's fingerprint rather than a shortcoming: transaction and
operating balances deliver liquidity services, are costly to relocate,
and have no natural decision point at which a depositor reconsiders the
bank; certificates mature, and the disclosure arrives while the renewal
decision is live.  The null also bounds the aggregate
interpretation---the CECL disclosure reallocated term funding at the
margin; it did not trigger broad deposit flight.

\subsection{Additional Checks}
\label{sec:robustness:additional}

Four further sets of results, reported compactly here (tables available
on request).  \emph{Treatment specification.}  Scaling the Day-One
charge by total assets or by loans instead of equity preserves the
pattern (uninsured time deposits $-0.714$, $p<0.10$; interest expense
and ROA significant at 1\%), and binary cutoffs at the 75th and 90th
percentiles deliver estimates of $-0.28$ and $-0.29$, so neither the
scaling nor the threshold is doing the work.  Dropping the single bank
with the largest adjustment leaves the baseline essentially unchanged.
\emph{Residualized treatment.}  Replacing the adjustment with the
residual from the expanded 2022Q4 predictor regression of
Table~\ref{tab:cecl_predictors_cross_section} column (2) preserves sign
at roughly two-thirds the baseline magnitude, though imprecisely,
suggesting the estimate is not driven purely by selection on
observables.  \emph{Local reallocation.}  In annual branch-level SOD
data, deposit growth after 2023 is faster at low-CECL banks more
exposed to high-CECL competitors in their counties ($+2.4$ percentage
points per unit of leave-out exposure, $p<0.10$; county-level analog
$+5.2$, $p<0.01$), so part of the withdrawn funding re-intermediates
locally---spillovers that bias the within-market estimates toward zero.
\emph{Lending.}  Quarterly loan growth at high-CECL banks is lower
after adoption relative to the four preceding quarters ($-0.65$,
$p<0.01$, concentrated in CRE), but growth in those four quarters was
unusually strong, and measured against a longer pre-period the slowdown
is modest at best; we therefore treat the lending margin as suggestive
only and draw no headline conclusion from it.
